
\section{Kết luận và hướng phát triển}

\subsection{Tóm tắt quá trình thực hiện}

Đề tài được thực hiện qua các bước: chuẩn bị và tiền xử lý
dữ liệu MNIST cùng dữ liệu chữ viết tay người Việt được nhóm thu thập; xây dựng,
huấn luyện và đánh giá bốn mô hình Multinomial Naive Bayes,
Logistic Regression, MLP và Random Forest; thử nghiệm bổ sung
dữ liệu đúng miền và tăng cường dữ liệu; cuối cùng xây dựng
giao diện demo bằng Gradio để trực quan hóa kết quả dự đoán
và so sánh thời gian xử lý giữa các mô hình.

\subsection{Kết luận}

Trong bốn mô hình được thử nghiệm, MLP đạt kết quả tốt nhất
trên tập kiểm thử MNIST với accuracy $98.09\%$. Khi áp dụng
trên chữ viết tay người Việt, accuracy của MLP giảm từ
$98.13\%$ xuống $89.35\%$. Việc bổ sung dữ liệu đúng miền
giúp accuracy tăng lên $94.56\%$ đối với MLP và từ $68.78\%$
lên $87.89\%$ đối với Random Forest.

Ngoài ra, tăng cường dữ liệu giúp MLP đạt accuracy $77.91\%$
khi ảnh bị dịch $3$ pixel, cao hơn đáng kể so với $20.19\%$
của mô hình không tăng cường dữ liệu.

\subsection{Hướng phát triển}

Trong tương lai, nhóm sẽ thử nghiệm mô hình CNN để khai
thác đặc trưng không gian của ảnh tốt hơn, đồng thời thu thập
thêm dữ liệu chữ viết tay tiếng Việt đa dạng về người viết
và kiểu chữ nhằm cải thiện khả năng tổng quát hóa của mô hình.
